\begin{proposition}[Finite field embedding fingerprints: the cubic characteristic polynomial of the second-order collision kernel manifests in higher-moment recursion (auditable)]\label{prop:pom-charpoly-modp-a2-embedding}
Denote the characteristic cubic of the second-order collision kernel \(A_2\) (Proposition \ref{prop:pom-collision-det}) by
\[
\chi_{A_2}(\lambda):=\lambda^3-2\lambda^2-2\lambda+2.
\]
On the minimal integer-coefficient recursion characteristic polynomials \(P_q(\lambda)\) within the resonance window (Theorem \ref{thm:pom-moment-resonance}), there exists an auditable finite field ``embedding'' phenomenon:
\[
\boxed{\;
\chi_{A_2}(\lambda)\mid P_{14}(\lambda)\ \ (\mathrm{mod}\ 7),
\qquad
\chi_{A_2}(\lambda)\mid P_{12}(\lambda)\ \ (\mathrm{mod}\ 11).\;}
\]
The corresponding explicit factorizations are given in Table \ref{tab:fold_collision_charpoly_mod7_selection_q13_15} and Table \ref{tab:fold_collision_charpoly_mod11_fingerprint_q12_15}.
The above finite field factorizations are generated by the script \texttt{scripts/exp\_fold\_collision\_charpoly\_modp\_factorization.py} directly from the \(P_q\) coefficients listed in Theorem \ref{thm:pom-moment-resonance} via reduction and factorization, requiring no new DP data.
\end{proposition}

\begin{remark}[Selection law modulo \texorpdfstring{$7$}{7}: a one-glance boundary between short in-field periods and long extension-field periods (auditable)]\label{rem:pom-charpoly-mod7-selection}
From the \(\mathbb{F}_7[\lambda]\) factorization in Table \ref{tab:fold_collision_charpoly_mod7_selection_q13_15}:
\begin{itemize}
  \item For \(q=13\), all nonzero eigenphases lie in \(\mathbb{F}_7^\ast\) (maximum irreducible degree is \(1\)), and repeated roots appear (e.g., \((\lambda-2)^3\)).
  Consequently, after aligning the finite transient (from \(\lambda^2\)), the recursion dynamics modulo \(7\) is forced into an auditable class of ``short period + polynomial modulation''; the crude upper bound read from the table is
  \[
  \mathrm{Per}_7(q=13)\ \mid\ 42.
  \]
  \item For \(q=15\), irreducible quadratic and irreducible cubic factors appear (maximum irreducible degree is \(3\)), so the phases must enter the extension fields \(\mathbb{F}_{7^2}\) and \(\mathbb{F}_{7^3}\);
  simultaneously, the squared factors \((\lambda\pm 1)^2\) force the presence of generalized eigenvalue modulation.
  Under the same crude-bound calibration, the unified upper bound read from the table is
  \[
  \mathrm{Per}_7(q=15)\ \mid\ 19152.
  \]
\end{itemize}
The above conclusions are read entirely from the finite field factorization of \(P_q(\lambda)\), with no dependence on new DP computations; they split the odd-order \(q\) within the same resonance window into two directly auditable classes of dynamical complexity under reduction modulo \(7\).
\end{remark}

\begin{remark}[Squared factors \texorpdfstring{$\Rightarrow$}{=>} polynomial modulation: an audit handle for recursion modulo \texorpdfstring{$11$}{11}]\label{rem:pom-charpoly-mod11-modulation}
Table \ref{tab:fold_collision_charpoly_mod11_fingerprint_q12_15} shows that in the modulo \(11\) slice, \(q=12\) exhibits a squared quadratic factor \((\lambda^2\pm 4\lambda-3)^2\), while \(q=13,15\) exhibit squared linear factors \((\lambda\pm 1)^2\) (with \(\mathrm{repeat?}=1\) for all three).
This implies that the corresponding finite field recursion operators possess nontrivial Jordan blocks under this prime shadow, so that their solution spaces contain ``polynomially modulated exponential phase'' (generalized eigenmode) components.
\par
Therefore, an auditable consistency check modulo \(11\) is the following: after reducing the recursion to \(\mathbb{F}_{11}\), if the observed sequence behavior is hypothesized to be ``purely periodic given only by semisimple phases'' (i.e., requiring no \(\times 11\) modulation factor), then one must further verify whether the minimal annihilating polynomial of the specific sequence over \(\mathbb{F}_{11}\) has already eliminated the squared factors (e.g., by refitting the finite field Hankel rank/minimal recursion once, which suffices for audit).
\end{remark}

\begin{remark}[Low-complexity fingerprint modulo \texorpdfstring{$13$}{13}: pure period upper bound for \texorpdfstring{$q=13$}{q=13}]\label{rem:pom-charpoly-mod13-small-period}
Table \ref{tab:fold_collision_charpoly_mod13_fingerprint_q13} gives: under reduction modulo \(13\), the nonzero irreducible factors of \(P_{13}(\lambda)\) have degree at most \(2\) (maximum irreducible degree is \(2\)), and apart from the transient factor \(\lambda^2\), no squared factors appear (\(\mathrm{repeat?}=0\)).
Therefore, after aligning the finite transient, the recursion dynamics modulo \(13\) belongs to the purely periodic class, with a crude period upper bound
\[
\mathrm{Per}_{13}(q=13)\ \mid\ 13^2-1=168,
\]
consistent with the upper bound read from the table.
\end{remark}

\begin{remark}[Cross-prime audit method: automatically generating modulo \texorpdfstring{$p$}{p} recursion fingerprints from \texorpdfstring{$P_q(\lambda)\bmod p$}{Pq mod p}]\label{rem:pom-charpoly-modp-audit-method}
For any given \((q,p)\), factor \(P_q(\lambda)\) over \(\mathbb{F}_p[\lambda]\) as
\(
P_q(\lambda)\equiv \prod_i f_i(\lambda)^{e_i}
 \).
This factorization directly yields a set of auditable fingerprints:
\begin{itemize}
  \item Whether squared factors appear (some \(e_i\ge 2\) with \(f_i\neq \lambda\)): this determines whether the ``generalized phase/polynomial modulation'' channel exists;
  \item The maximum irreducible degree \(\max_i \deg f_i\): this gives the minimal extension field degree required, thereby providing an upper bound factor \((p^{\deg f_i}-1)\) on the phase order;
  \item If all nonzero factors are linear: then the phase set lies entirely in \(\mathbb{F}_p^\ast\), and a minimal period upper bound can be directly given (multiplied by a possible modulation factor \(p\)).
\end{itemize}
Compared with an audit using only the recursion order \(d_q\), this ``finite field factorization fingerprint'' simultaneously provides explanatory power (where extension field modes arise, where Jordan modulation occurs) and computable upper bounds (e.g., the period bound column in Tables \ref{tab:fold_collision_charpoly_mod7_selection_q13_15}--\ref{tab:fold_collision_charpoly_mod13_fingerprint_q13}).
\end{remark}

\input{sections/generated/tab_fold_collision_charpoly_mod7_selection_q13_15_en}
\input{sections/generated/tab_fold_collision_charpoly_mod11_fingerprint_q12_15_en}
\input{sections/generated/tab_fold_collision_charpoly_mod13_fingerprint_q13_en}

